%! AP Statistics — Unit 2, Lesson 2.9 — Guided Notes
\documentclass[10pt]{article}
\usepackage{apstats-article}
\usepackage{apstats-boxes}

\begin{document}

\pageheader{Unit 2, Lesson 2.9}{Guided Notes}
\namedateperiod

\vspace{0.12in}

% ── OBJECTIVE ────────────────────────────────────────────────────────────────
\begin{objectivebox}
\small By the end of this lesson, I will be able to\dots\\[4pt]
\writeline\\[1pt]
\writeline\\[1pt]
\writeline\\[1pt]
\writeline
\end{objectivebox}

\vspace{0.1in}

% ── VOCABULARY ───────────────────────────────────────────────────────────────
\begin{vocabbox}
\termblanklong{Nonlinear relationship}
\termblanklong{Transformation}
\termblanklong{Exponential model ($y = a \cdot b^x$)}
\termblanklong{Power model ($y = a\,x^b$)}
\termblanklong{Back-transformation}
\termblanklong{Linearization check}
\end{vocabbox}

\vspace{0.1in}

% ── HOOK ─────────────────────────────────────────────────────────────────────
\begin{hookbox}
\small\textbf{Opening Question:} A scatterplot of bacteria count (millions) vs.\ hours of
growth curves upward sharply. A student fits a linear regression and gets a curved residual
plot. What does the residual plot tell us? What should we do next?

\vspace{6pt}
What the residual plot tells us: \blank{7cm}

\vspace{4pt}
\textbf{Key idea:} Instead of abandoning regression, we can \textbf{transform} the data to
make the relationship \blank{3cm}, then apply the linear regression tools we already know.
\end{hookbox}

\vspace{0.1in}

% ── STEP 1 ───────────────────────────────────────────────────────────────────
\begin{notesbox}{Step 1 --- Diagnosing Nonlinearity (DAT-1.G)}
\small

\noindent A linear model is \textbf{not appropriate} when:
\begin{itemize}[itemsep=2pt]
  \item The \textbf{scatterplot} shows a \blank{4cm} pattern (not a straight line), \textbf{or}
  \item The \textbf{residual plot} shows a \blank{4cm} pattern (not random scatter).
\end{itemize}

\vspace{6pt}
\textbf{Example:} Bacteria count (millions) $y$ vs.\ hours $x$. \\[4pt]
\noindent Original scatterplot: \blank{5cm}. \quad Original residual plot: \blank{5cm}.

\vspace{6pt}
\textbf{Conclusion:} The linear model is \blank{3cm} appropriate because \blank{6cm}.

\end{notesbox}

\vspace{0.1in}

% ── STEP 2 ───────────────────────────────────────────────────────────────────
\begin{notesbox}{Step 2 --- The Exponential Model: Transform $y$ (DAT-1.H)}
\small

\noindent \textbf{When to use:} The scatterplot looks like an exponential curve
($y$ grows or decays multiplicatively).

\vspace{6pt}
\textbf{Transformation:} Apply $\log_{10}$ to the \blank{2cm} values only.

\vspace{4pt}
\textbf{Why it works:}
\[
  y = a \cdot b^x
  \;\;\Longrightarrow\;\;
  \log y = \log a + (\log b)\,x
  \;\;\Longrightarrow\;\;
  \underbrace{\log y}_{\text{new } y'} = \underbrace{\log a}_{A} + \underbrace{(\log b)}_{B}\,x
\]
This is the equation of a \blank{6cm} in $(\log y,\, x)$.

\vspace{8pt}
\textbf{Regression of $\log y$ on $x$} gives: \quad $\widehat{\log y} = $ \blank{6cm}

\vspace{6pt}
\textbf{Back-transform} to predict on the original scale:
\[
  \hat{y} = 10^{\widehat{\log y}} = 10^{a + bx}
\]

\vspace{6pt}
\textbf{Example} --- bacteria count data. Technology gives: $\widehat{\log(\text{count})} = 0.85 + 0.22x$

\vspace{4pt}
Write the back-transformed exponential equation: \blank{7cm}

\vspace{4pt}
Predict count at $x = 6$ hours: \blank{7cm}

\end{notesbox}

\vspace{0.1in}

% ── STEP 3 ───────────────────────────────────────────────────────────────────
\begin{notesbox}{Step 3 --- The Power Model: Transform Both $x$ and $y$ (DAT-1.H)}
\small

\noindent \textbf{When to use:} The scatterplot curves like a power function ($y$ grows as a
power of $x$, e.g.\ area vs.\ radius).

\vspace{6pt}
\textbf{Transformation:} Apply $\log_{10}$ to \blank{5cm} values.

\vspace{4pt}
\textbf{Why it works:}
\[
  y = a\,x^b
  \;\;\Longrightarrow\;\;
  \log y = \log a + b\,\log x
\]

\vspace{6pt}
\textbf{Regression of $\log y$ on $\log x$} gives: \quad $\widehat{\log y} = $ \blank{6cm}

\vspace{6pt}
\textbf{Back-transform:} \quad $\hat{y} = 10^{\widehat{\log y}}$

\vspace{6pt}
\renewcommand{\arraystretch}{1.6}
\begin{tabularx}{\linewidth}{@{} >{\bfseries}p{3cm} X X @{}}
\rowcolor{sky} \textbf{Model} & \textbf{Transformation} & \textbf{Regression} \\
Exponential & $\log y$ only & $\log y$ on $x$ \\
Power       & $\log y$ and $\log x$ & $\log y$ on $\log x$ \\
\end{tabularx}

\end{notesbox}

\vspace{0.1in}

% ── STEP 4 ───────────────────────────────────────────────────────────────────
\begin{notesbox}{Step 4 --- Checking the Transformation (DAT-1.G)}
\small

\noindent After transforming, \textbf{always} inspect the \blank{4cm} of the transformed
model.

\vspace{4pt}
\begin{itemize}[itemsep=3pt]
  \item \textbf{Success:} residual plot shows \blank{5cm} --- the transformation worked.
  \item \textbf{Failure:} residual plot still shows a \blank{4cm} --- try a different transformation.
\end{itemize}

\vspace{6pt}
Also compare \textbf{$r^2$}: a good transformation should \blank{5cm} $r^2$ compared to
the original linear fit.

\end{notesbox}

\vspace{0.1in}

% ── CONNECTIONS ───────────────────────────────────────────────────────────────
\begin{spiralbox}
\small
\begin{multicols}{2}

\textbf{How this connects:}
\begin{itemize}[itemsep=3pt]
  \item Lessons 2.7--2.8: residual plots diagnose model fit; regression formulas apply
        to transformed data exactly as before.
  \item Unit 9: inference on slopes requires linearity --- transformations satisfy this
        condition for nonlinear data.
\end{itemize}

\columnbreak

\textbf{Reflection:}\\
A friend says: ``I transformed $y$ and the slope of the transformed regression is $b = 0.22$.
That means $y$ increases by 0.22 for each unit increase in $x$.'' What is wrong with this
interpretation?\\[2pt]
\writeline\\[1pt]\writeline\\[1pt]\writeline

\end{multicols}
\end{spiralbox}

\vspace{0.1in}

\begin{notesbox}{Additional Notes}
\vspace{0pt}
\writeline\\\writeline\\\writeline\\\writeline\\\writeline\\\writeline
\end{notesbox}

\end{document}
